\section{Delta Wing Mesh Validation: OpenFOAM vs Wind Tunnel vs ANSYS Fluent}


To verify that the CFD results are not strongly affected by mesh resolution, two meshes with different refinement levels were generated and compared. The first mesh (Mesh~1) contains approximately $8.74\times 10^{6}$ cells with a target $y^{+}\approx 4$, while the second mesh (Mesh~2) contains approximately $1.93\times 10^{6}$ cells with a target $y^{+}\approx 5$. Both meshes were generated with 10 prism layers on the aircraft surface. The mesh comparison was performed using the Spalart--Allmaras (SA) turbulence model, and both meshes were simulated over the full angle of attack range from $-4^{\circ}$ to $30^{\circ}$. The results are summarised in Table~\ref{tab:mesh_SA}.

At low and moderate angles of attack (from $-4^{\circ}$ to approximately $10^{\circ}$), both meshes provide very similar aerodynamic coefficients. In this range, the lift coefficient differs by less than about 1\% between the meshes. For example, at $10^{\circ}$, Mesh~1 predicts $C_{L}=0.5635$ while Mesh~2 predicts $C_{L}=0.5684$, which corresponds to a difference below 0.9\%. The drag coefficient is also nearly identical in this range, with differences typically below 1\%. These results indicate that both meshes provide sufficient resolution when the flow is mostly attached.

At higher angles of attack (above approximately $14^{\circ}$), the differences between the meshes become more visible. The coarser Mesh~2 tends to predict slightly higher lift and drag compared to the finer Mesh~1. At $20^{\circ}$, Mesh~1 predicts $C_{L}=1.0199$ and $C_{D}=0.3822$, while Mesh~2 predicts $C_{L}=1.0722$ and $C_{D}=0.3970$, which corresponds to a difference of roughly 5\% in lift and 4\% in drag. At $30^{\circ}$, the lift values become close again ($C_{L}=1.3211$ for Mesh~1 and $C_{L}=1.3116$ for Mesh~2), and the drag values are also very similar ($C_{D}=0.7599$ for Mesh~1 and $C_{D}=0.7534$ for Mesh~2).

The pitching moment coefficient shows a higher sensitivity to the mesh at high angles of attack, which is expected because the pitching moment depends on the detailed pressure distribution over the full surface and is therefore more affected by local mesh differences. For instance, at $22^{\circ}$, Mesh~1 predicts $C_{My,cog}=0.0626$ while Mesh~2 predicts $C_{My,cog}=0.0820$. Even though the absolute values differ more than for lift and drag, the overall trend of the pitching moment curve remains consistent.

When compared to the wind tunnel data and the ANSYS Fluent results reported by Giannelis et al.~\cite{Giannelis2023SSAMGen5}, both meshes show good agreement in the low to moderate angle of attack range. At high angles of attack, some deviation is expected because steady RANS simulations have known limitations for flows with strong separation, vortex interaction, and possible vortex breakdown. Similar deviations were also reported in the ANSYS Fluent results in the same reference.

Based on these comparisons, Mesh~1 (approximately $8.7$ million cells) was selected for all subsequent simulations. Although Mesh~2 provides reasonable predictions across most of the angle of attack range, the finer mesh gives more consistent results at high angles of attack, especially for the pitching moment coefficient.

\section{Delta Wing Turbulence Model Validation: OpenFOAM vs Wind Tunnel vs ANSYS Fluent}


After selecting the finer mesh (Mesh~1, $8.74\times 10^{6}$ cells, $y^{+}\approx 4$), the influence of the turbulence model was evaluated by comparing the Spalart--Allmaras (SA) and $k$--$\omega$ SST models. The results are summarised in Table~\ref{tab:tm_mesh1}.

At low angles of attack (from $-4^{\circ}$ to approximately $8^{\circ}$), both turbulence models produce very similar results. For example, at $6^{\circ}$, the SA model predicts $C_{L}=0.3375$ and $C_{D}=0.0573$, while the $k$--$\omega$ SST model predicts $C_{L}=0.3284$ and $C_{D}=0.0560$. In this range the flow remains mostly attached, so the solution is less sensitive to the turbulence model choice.

From approximately $10^{\circ}$ and higher, the differences between the models become more visible. In the mid-range between $14^{\circ}$ and $20^{\circ}$, the $k$--$\omega$ SST model predicts slightly higher lift compared to SA (for example, at $20^{\circ}$: $C_{L}=1.0756$ for SST vs $C_{L}=1.0199$ for SA). The drag coefficient follows a similar tendency, with SST predicting slightly higher drag at many high angles.

The largest differences are observed in the pitching moment coefficient at higher angles of attack. This behaviour is expected because the pitching moment is highly sensitive to small changes in pressure distribution, which depend on how each turbulence model represents separation and vortex structures.

Overall, the trends obtained in OpenFOAM are consistent with the behaviour reported by Giannelis et al.~\cite{Giannelis2023SSAMGen5} for ANSYS Fluent, and both turbulence models reproduce the general shape of the lift, drag, and pitching moment curves across the full angle of attack range.

% ------------------------------------------------------------
% Tables
% ------------------------------------------------------------

\begin{table}[!ht]
\centering
\caption{Mesh comparison using the Spalart--Allmaras turbulence model: Mesh~1 ($8.74\times10^{6}$ cells, $y^{+}\approx 4$) vs Mesh~2 ($1.93\times10^{6}$ cells, $y^{+}\approx 5$). Both meshes use 10 prism layers.}
\label{tab:mesh_SA}
\setlength{\tabcolsep}{6pt}
\renewcommand{\arraystretch}{1.15}
\begin{tabular}{c c c c | c c c c}
\hline
\multicolumn{4}{c|}{\textbf{SA --- Mesh 1: $8{,}741{,}586$ cells, $y^{+}\approx 4$, 10 layers}} &
\multicolumn{4}{c}{\textbf{SA --- Mesh 2: $1{,}926{,}908$ cells, $y^{+}\approx 5$, 10 layers}}\\
\hline
$\alpha$ [$^\circ$] & $C_L$ & $C_D$ & $C_{My,cog}$ &
$\alpha$ [$^\circ$] & $C_L$ & $C_D$ & $C_{My,cog}$\\
\hline
-4 & -0.1978 & 0.0362 & -0.0131 & -4 & -0.2343 & 0.0426 & -0.0249 \\
-2 & -0.1181 & 0.0324 & -0.0125 & -2 & -0.1187 & 0.0328 & -0.0123 \\
0  & -0.0091 & 0.0292 & -0.0004 & 0  & -0.0091 & 0.0297 & -0.0006 \\
2  & 0.1007  & 0.0314 & 0.0109  & 2  & 0.1010  & 0.0321 & 0.0100 \\
4  & 0.2144  & 0.0397 & 0.0244  & 4  & 0.2145  & 0.0399 & 0.0232 \\
6  & 0.3375  & 0.0573 & 0.0388  & 6  & 0.3398  & 0.0575 & 0.0365 \\
8  & 0.4570  & 0.0836 & 0.0552  & 8  & 0.4599  & 0.0837 & 0.0529 \\
10 & 0.5635  & 0.1170 & 0.0732  & 10 & 0.5684  & 0.1170 & 0.0729 \\
12 & 0.6743  & 0.1592 & 0.0856  & 12 & 0.6797  & 0.1591 & 0.0834 \\
14 & 0.7901  & 0.2101 & 0.0884  & 14 & 0.8019  & 0.2116 & 0.0854 \\
16 & 0.8828  & 0.2653 & 0.0826  & 16 & 0.9110  & 0.2705 & 0.0843 \\
18 & 0.9679  & 0.3256 & 0.0779  & 18 & 1.0063  & 0.3344 & 0.0862 \\
20 & 1.0199  & 0.3822 & 0.0753  & 20 & 1.0722  & 0.3970 & 0.0899 \\
22 & 1.0853  & 0.4502 & 0.0626  & 22 & 1.1121  & 0.4568 & 0.0820 \\
24 & 1.1309  & 0.5119 & 0.0943  & 24 & 1.1636  & 0.5247 & 0.0781 \\
26 & 1.2046  & 0.5903 & 0.0949  & 26 & 1.2354  & 0.6048 & 0.0784 \\
28 & 1.2489  & 0.6651 & 0.1032  & 28 & 1.2845  & 0.6817 & 0.0822 \\
30 & 1.3211  & 0.7599 & 0.0866  & 30 & 1.3116  & 0.7534 & 0.0844 \\
\hline
\end{tabular}
\end{table}

\begin{table}[!ht]
\centering
\caption{Turbulence model comparison on Mesh~1 ($8.74\times10^{6}$ cells, $y^{+}\approx 4$): Spalart--Allmaras vs $k$--$\omega$ SST.}
\label{tab:tm_mesh1}
\setlength{\tabcolsep}{6pt}
\renewcommand{\arraystretch}{1.15}
\begin{tabular}{c c c c | c c c c}
\hline
\multicolumn{4}{c|}{\textbf{Mesh 1 --- SA}} &
\multicolumn{4}{c}{\textbf{Mesh 1 --- $k$--$\omega$ SST}}\\
\hline
$\alpha$ [$^\circ$] & $C_L$ & $C_D$ & $C_{My,cog}$ &
$\alpha$ [$^\circ$] & $C_L$ & $C_D$ & $C_{My,cog}$\\
\hline
-4 & -0.1978 & 0.0362 & -0.0131 & -4 & -0.2084 & 0.0375 & -0.0081 \\
-2 & -0.1181 & 0.0324 & -0.0125 & -2 & -0.0923 & 0.0292 & -0.0169 \\
0  & -0.0091 & 0.0292 & -0.0004 & 0  & -0.0055 & 0.0277 & 0.0034 \\
2  & 0.1007  & 0.0314 & 0.0109  & 2  & 0.0981  & 0.0281 & 0.0115 \\
4  & 0.2144  & 0.0397 & 0.0244  & 4  & 0.2061  & 0.0382 & 0.0297 \\
6  & 0.3375  & 0.0573 & 0.0388  & 6  & 0.3284  & 0.0560 & 0.0437 \\
8  & 0.4570  & 0.0836 & 0.0552  & 8  & 0.4395  & 0.0814 & 0.0642 \\
10 & 0.5635  & 0.1170 & 0.0732  & 10 & 0.5500  & 0.1149 & 0.0803 \\
12 & 0.6743  & 0.1592 & 0.0856  & 12 & 0.6483  & 0.1536 & 0.0903 \\
14 & 0.7901  & 0.2101 & 0.0884  & 14 & 0.7643  & 0.2053 & 0.0893 \\
16 & 0.8828  & 0.2653 & 0.0826  & 16 & 0.8942  & 0.2667 & 0.0915 \\
18 & 0.9679  & 0.3256 & 0.0779  & 18 & 1.0012  & 0.3320 & 0.0914 \\
20 & 1.0199  & 0.3822 & 0.0753  & 20 & 1.0756  & 0.3980 & 0.1004 \\
22 & 1.0853  & 0.4502 & 0.0626  & 22 & 1.1426  & 0.4684 & 0.1069 \\
24 & 1.1309  & 0.5119 & 0.0943  & 24 & 1.1951  & 0.5366 & 0.0938 \\
26 & 1.2046  & 0.5903 & 0.0949  & 26 & 1.2208  & 0.6002 & 0.0819 \\
28 & 1.2489  & 0.6651 & 0.1032  & 28 & 1.2604  & 0.6720 & 0.0893 \\
30 & 1.3211  & 0.7599 & 0.0866  & 30 & 1.3176  & 0.7634 & 0.0741 \\
\hline
\end{tabular}
\end{table}

